%% Determines the type of document (including standard settings for layout).
\documentclass[letterpaper]{article}



%% Package to control the size of a page: the standard margins used by LaTeX are
%% very wide!
\usepackage[margin=1in]{geometry}


%% Packages add functionality to the document. The AMS packages are standard
%% packages to support various mathematical notations. AMS stands for American
%% Mathematical Society, the organization that maintains these packages.
\usepackage{amsmath,amsthm,amssymb}

%% Theorems-like environments using functionality provided by amsthm.
\theoremstyle{plain}
\newtheorem{theorem}{Theorem}[section]
    %% [section] at the end specifies that theorems should be numbered
    %% per-section: Section x starts with theorem-like x.1 and so on, ...
\newtheorem{proposition}[theorem]{Proposition}
    %% [theorem] in the middle specifies: use the same counter as the theorem
    %% environment: here we number all theorem-like environments consecutively.
\newtheorem{corollary}[theorem]{Corollary}
\newtheorem{lemma}[theorem]{Lemma}
\theoremstyle{definition}
\newtheorem{definition}[theorem]{Definition}
\theoremstyle{remark}
\newtheorem{example}[theorem]{Example}
\newtheorem{remark}[theorem]{Remark}


%% Support for nicely formatted tables.
\usepackage{booktabs}


%% Support for colors & colors in tables.
\usepackage[table]{xcolor}

% Seven colors safe for use color blindness.
% Colors taken from doi:10.1038/nmeth.1618.
\definecolor{cbOrange}{RGB}{230,159,0}
\definecolor{cbSkyBlue}{RGB}{86,180,233}
\definecolor{cbBluishGreen}{RGB}{0,158,115}
\definecolor{cbYellow}{RGB}{240,228,66}
\definecolor{cbBlue}{RGB}{0,114,178}
\definecolor{cbVermillion}{RGB}{213,94,0}
\definecolor{cbReddischPurple}{RGB}{204,121,167}


%% Notation used in this document.
\newcommand{\n}{\mathbf{n}} %% Num. Replicas.
\newcommand{\f}{\mathbf{f}} %% Num. Faulty Replicas.

%% Misc. Math notation.
\newcommand{\BigO}{\mathcal{O}}
\newcommand{\abs}[1]{\lvert #1 \rvert}
\newcommand{\AName}[1]{\textsc{#1}}
\newcommand{\Var}[1]{\texttt{#1}}


%% Algorithms.
\usepackage{algorithm}
\usepackage[noend]{algorithmic}
\newcommand{\GETS}{:=}


%% Formatting SI-units.
\usepackage{siunitx}
\sisetup{per-mode=symbol}


%% TikZ: for creating figures.
\usepackage{tikz}

%% Configuration for figures: Nicer arrows.
\usetikzlibrary{arrows.meta}
\tikzset{>=Stealth}


%% pgfplots: drawing plots using TikZ.
\usepackage{pgfplots}
%% Configuration for plots: Use color-blind friendly colors.
\pgfplotscreateplotcyclelist{cbSafeList}{
    very thick,solid,cbOrange,every mark/.append style={solid},mark=*\\
    very thick,solid,cbSkyBlue,every mark/.append style={solid},mark=*\\
    very thick,solid,cbBluishGreen,every mark/.append style={solid},mark=*\\
    very thick,solid,cbYellow,every mark/.append style={solid},mark=*\\
    very thick,solid,cbBlue,every mark/.append style={solid},mark=*\\
    very thick,solid,cbVermillion,every mark/.append style={solid},mark=*\\
    very thick,solid,cbReddischPurple,every mark/.append style={solid},mark=*\\
    very thick,solid,black,every mark/.append style={solid},mark=*\\
}
\pgfplotsset{
    legend style={font=\small},
    compat=1.16,
    width=260pt,
    height=140pt,
    legend cell align=left,
    xlabel near ticks,
    ylabel near ticks,
    every axis/.append style={
        cycle list name=cbSafeList,
        ymin=0,
        enlargelimits=0.05,
        mark size=1pt,
        ylabel style={align=center},
        xlabel style={align=center},
        title style={align=center}
    }
}

%% PgfplotsTable: loading data files to use with pgfplots.
\usepackage{pgfplotstable}


%% Support for hyperlinks and urls. The setting ``colorlinks'' sets how links
%% are shown in the document (with a color, without underline). We put hyperref
%% last---it has a tendency to break other packages when loaded before them.
\usepackage[colorlinks]{hyperref}
\hypersetup{
    linkcolor=blue,
    citecolor=blue,
    urlcolor=blue
}
\usepackage{tocloft}
% Make Table of Contents and List of Figures text blue
\renewcommand{\cftsecfont}{\color{blue}}
\renewcommand{\cftsecpagefont}{\color{blue}}

\renewcommand{\cftfigfont}{\color{blue}}
\renewcommand{\cftfigpagefont}{\color{blue}}
\usepackage{graphicx}
\usepackage{pgffor}
\usepackage{caption}
\usepackage{tabularx}
\usepackage{enumitem}
\usepackage{fancyhdr}
\usepackage{float}

%% Page numbering
\pagestyle{fancy}
\fancyhead[R]{\textit{Page \thepage}}
\fancyhead[L]{\textit{LAB REPORT 3}}
\fancyfoot[C]{}
\renewcommand{\headrulewidth}{0pt}
\renewcommand{\footrulewidth}{0pt}

\newcounter{experiment}

\begin{document}

\begin{titlepage}

	\thispagestyle{empty}
	\centering
	\vspace*{3em}

	\Huge{COMPSCI 2XC3 Final Lab Report}\\[1em]
	\Large Prepared by\\[2em]
	\Large Group 64\\[2em]

	\LARGE Luca Mawyin\\[0.5em]
	\LARGE Anderson Ray\\[0.5em]
	\LARGE Theodore Pham\\[2em]

	\Large COMPSCI 2ME3 \\[0.5em]
	\Large McMaster University\\[2em]
	\Large \today\\[2cm]

\end{titlepage}

\tableofcontents
\newpage

\listoffigures
\newpage

\listoftables
\newpage


\stepcounter{section}
\section*{Part \thesection}
\addcontentsline{toc}{section}{Part \thesection}

\stepcounter{experiment}
\subsection*{Experiment \theexperiment}
\label{exp1}
\addcontentsline{toc}{subsection}{Experiment \theexperiment}

\subsubsection*{Objective}
\subsubsection*{Results}
\subsubsection*{Discussion}

\stepcounter{section}
\section*{Part \thesection}
\addcontentsline{toc}{section}{Part \thesection}

The A* Algorithm is an algorithm used find the shortest path from a start node to a target node in a directed weighted graph. It's algorithm is practically the same as Dijkstra's, the only difference is how it determines priority (the key value) in the priority queue.
\\

In A*, the key value used in the priority queue is equal to the total weight of the edges along the path traversed (like Dijkstra's) plus some value coming from a heuristic function which is a lower bound to the weight of the rest of the path to the target.

A* does this to address a problem that Dijkstra has where it explores nodes that are obviously not on the shortest path. For example if I were to calculate the shortest route from Toronto to Hamilton with Dijktra's it'll probablys venture out to Scarborough before finding Hamilton which is obviously wrong. If I were to use A* though, I could use a heurisitc function that calculates the straight line distance to Hamilton. This heurisitc function would make A* prioritize exploring western locations (relative to Toronto) aside from exploring Scarborough. This would result in the shortest path being found much quicker.
\\

You could empirically test the speed difference between A* to Dijkstra's by generating a varying number of points in 2D space, with random edges going between the two and running the algorithm on random starting and target nodes, comparing the execution times. It is important to note that for any node, it's heuristic value (straight line distance in this case) must be a lower bound to the sum of the weights of the path to the target node, otherwise A* could break and return a path that isn't the shortest.
\\

If you generated an arbitrary heuristic function, A* could return the wrong result (same reason mentioned prior) and could also take longer then Dijktra's.
\\

Applications where A* should be used is when there exists a heurstic function that is quick to compute, is always a lower bound to the weight of the shortest path from the inputted node to the target node, and the function decreases as you move to nodes closer to the target. This is why A* is so good for shortest route calculations because straight line distances are always the lower bound to the shortest route, straight line distances are easy to calculate ($O(1)$ via pythagorean theorem if you have cartesian coordinates) and it decreases as you get closer to the target node.    



% A* is used in apps such as google maps for calculating the shortest route to a destination. In scenario's like shortest route calculators, the nodes of the graph would be intersections with their distances between them being the weights of the edges between them, and the heuristic function would return the straight line distance between the inputted node to the targeted node.


\end{document}
